% ~3 min
\begin{frame*}[vertical-alignment = top]
  \frametitle{A Specification and a Program}
  \vspace*{6mm}
  \[ \op{map}\,(\lambda x.\, 2x) \circ \op{sort} \]
  \ocamlblock{sort-times-two.ml}
  \begin{uncoverenv}<2->
    \begin{center}
      Can we \emph{derive} the program from the specification?
    \end{center}
  \end{uncoverenv}
  \speakernote{Top: sort, then double every element (two traversals).
    Bottom: one traversal that doubles each pivot as it goes. They are
    equal, but proving it by hand is fiddly. Goal of the talk: derive the
    second from the first, every step checked by Rocq.}
\end{frame*}

\begin{frame}
  \frametitle{Program Calculation}
  \begin{center}
    \begin{tikzpicture}[alt={A simple specification is rewritten by
        algebraic laws, step by step, into an efficient program; each step
        is an equation.},
      stage/.style = {text width = 26mm, align = center,
        minimum height = 12mm}]
      \node[box, stage] (s) {simple\\ specification};
      \node[box, stage, right = 14mm of s] (m) {intermediate\\ program};
      \node[hot, stage, right = 14mm of m] (p) {efficient\\ program};
      \draw[->, thick] (s) -- node[above, note] {$=$} (m);
      \draw[->, thick] (m) -- node[above, note] {$=$} (p);
    \end{tikzpicture}
  \end{center}
  \begin{itemize}[<+->]
    \item Every step is an equation: correct by construction.
    \item The laws come from \embf{recursion schemes}.
  \end{itemize}
  \speakernote{Equational reasoning in the style of Bird and Meertens
    (Algebra of Programming); GHC's fusion rules are the same idea. Every
    arrow is an equation, so the final program equals the specification.
    The laws we need come from recursion schemes: next part.}
\end{frame}
